% ========== PREÂMBULO DO DOCUMENTO ==========
% Define a classe de documento como article (artigo)
\documentclass[12pt]{article}

% ========== PACOTES NECESSÁRIOS ==========
% Pacote para configurar a codificação de entrada como UTF-8
\usepackage[utf8]{inputenc}

% Pacote para configurar a codificação de saída com suporte a caracteres acentuados
\usepackage[T1]{fontenc}

% Pacote para configurar margens: esquerda e direita = 2cm, topo = 1.5cm, fundo = 2.5cm
\usepackage[left=2cm, right=2cm, top=1.5cm, bottom=2.5cm]{geometry}

% Pacote para usar fonte Arial em todo o documento
\usepackage{helvet}
% Comando para definir Arial como fonte padrão
\renewcommand{\familydefault}{\sfdefault}

% Remove o recuo (indentação) do início dos parágrafos
\setlength{\parindent}{0pt}

% Define o espaçamento vertical entre parágrafos
\setlength{\parskip}{6pt}

% ========== DEFINIÇÃO DE COMANDOS PERSONALIZADOS ==========
% Comando para o número do protocolo (padrão: 0000)
\newcommand{\protocolo}{010.601.50225-31}
% Comando para o tipo de reclamação (padrão: RECLAMAÇÃO xxx)
\newcommand{\reclamacao}{REC 007\_2025 - SÃO JOSÉ DE ESPIRANHAS - PERDAS NOS REATORES}
% Comando para a data (padrão: Mês de ano)
\newcommand{\data}{maio de 2025}

% ========== INÍCIO DO DOCUMENTO ==========
\begin{document}

% Título da seção
\begin{center}\textbf{ALEGAÇÃO ADICIONAL PARA ANEXAR AO PROTOCOLO Nº \protocolo{}}\end{center}
\vspace{0.5cm}
Prezado ouvidor,

A Reclamação Administrativa referente a \textbf{\reclamacao{}} foi devidamente protocolada em \data{}. Contudo, verifica-se inconsistência nos documentos apresentados pela concessionária.

% ========== ESPAÇO PARA CONTESTAÇÃO ========== %
Em análise ao memorial de cálculo apresentado pela concessionária Energisa, verifica-se que houve erro na forma de aplicação da restituição em dobro referente às cobranças indevidas relacionadas às perdas nos reatores de iluminação pública.

Embora a concessionária tenha reconhecido e apurado valores cobrados indevidamente ao longo de todo o período de 10 (dez) anos, a devolução em dobro foi aplicada apenas aos valores referentes ao período de 05/2015 a 04/2020, desconsiderando os valores apurados entre 05/2020 e 04/2025.

Dessa forma, o cálculo apresentado resulta em restituição parcial e incorreta, pois o correto é que a devolução em dobro incida sobre a totalidade do valor apurado nos 10 (dez) anos reconhecidos pela própria concessionária.

Assim, requer-se a revisão do memorial de cálculo e a retificação dos valores, com a aplicação da restituição em dobro sobre todo o período considerado na apuração da cobrança indevida.
% ========== FIM DO DOCUMENTO ==========
\end{document}
